%=====================================================================
\section{The Unique Dimension: Why $d = 5$}
\label{ch:whyd5}
%=====================================================================

Chapter~\ref{ch:whyC} established that points carry vectors $\psi_i \in \mathbb{C}^d$ for some $d$. We now prove that $d = 5$ is the unique value producing non-trivial physics. The argument proceeds in four steps: (i) identify the ``atomic'' dimensions, (ii) show that repeated atoms annihilate via swap symmetry, (iii) show that non-atomic dimensions decompose into atoms, and (iv) verify that only $d = 2 + 3 = 5$ survives.

\subsection{Atomic dimensions}

\begin{definition}[Atomic dimension]
An integer $n \geq 2$ is \emph{atomic} if it cannot be written as $n = a + b$ with $a, b \geq 2$. The integer $n = 1$ is excluded: a 1-dimensional complex vector space $\mathbb{C}^1$ supports only $\text{U}(1)$ (pure phase), which is not an independent gauge sector but the relative phase between sectors (cf.\ Sec.~\ref{sec:glue}).
\end{definition}

\begin{lemma}[The atoms are $\{2, 3\}$]
\label{lem:atoms}
The only atomic dimensions are $2$ and $3$.
\end{lemma}

\begin{proof}
For $n \geq 4$: write $n = 2 + (n - 2)$, where $n - 2 \geq 2$. Hence $n$ is not atomic.

For $n = 2$: the only partition into parts $\geq 1$ is $2 = 1 + 1$, but $1 < 2$. No valid decomposition exists. Atomic.

For $n = 3$: the partitions are $3 = 1 + 2$ and $3 = 1 + 1 + 1$. Since $1 < 2$, no valid decomposition exists. Atomic.
\end{proof}

\begin{remark}
Every integer $n \geq 2$ can be written as $n = 2a + 3b$ for non-negative integers $a, b$. The atoms $\{2, 3\}$ are the additive generators of all dimensions above the ``phase-only'' threshold.
\end{remark}

\subsection{Swap elimination: Group-theoretic proof}

\begin{theorem}[Swap Automorphism]
\label{thm:swap}
Let $\mathbb{C}^d = \mathbb{C}^a \oplus \mathbb{C}^a$ with $a \geq 2$. The gauge group $G = \text{SU}(a) \times \text{SU}(a) \times \text{U}(1)$ admits an outer automorphism $\sigma$ that forces the vacuum to be symmetric. Symmetric vacua produce no physical structure.
\end{theorem}

\begin{proof}
\textbf{Step 1: The automorphism.}
Define the swap:
\begin{equation}
\sigma: (g_1, g_2, e^{i\theta}) \;\mapsto\; (g_2, g_1, e^{-i\theta}).
\end{equation}
Since $\text{SU}(a)_1 \cong \text{SU}(a)_2$, this is a well-defined outer automorphism with $\sigma^2 = \text{id}$, generating $\mathbb{Z}_2$.

\textbf{Step 2: Vacuum classification.}
Let $|\Omega\rangle$ be a vacuum state.

\emph{Case (a): $\sigma|\Omega\rangle = |\Omega\rangle$.}
The two sectors are physically identical. They freeze at equal rates. By the DRLT simulation result (EXP-029): symmetric freezing preserves $\hbar_\text{eff}$ and all $W$ ratios. No coupling splitting, no force hierarchy, no structure.

\emph{Case (b): $\sigma|\Omega\rangle \neq |\Omega\rangle$.}
Spontaneous $\mathbb{Z}_2$ breaking. Domain walls form between $|\Omega\rangle$ and $\sigma|\Omega\rangle$. Tunneling restores symmetry in the ensemble average. The long-range physics is again symmetric.

In both cases, no net asymmetry survives. Physics is trivial.
\end{proof}

\begin{corollary}[Repeated atoms annihilate]
\label{cor:annihilate}
Any decomposition $d = n_1 + n_2 + \cdots + n_k$ containing a repeated pair $n_i = n_j$ admits a swap $\sigma_{ij}$ on that pair. By Theorem~\ref{thm:swap}, the pair annihilates. The effective dimension reduces by $2n_i$.
\end{corollary}

\begin{example}
$d = 4 = \mathbf{2+2}$: swap $\to$ trivial. \;
$d = 6 = \mathbf{3+3}$: swap $\to$ trivial. \;
$d = 7 = \mathbf{2+2}+3$: the 2-pair annihilates, leaving only 3. \;
$d = 8 = 2+\mathbf{3+3}$: the 3-pair annihilates, leaving only 2. \;
$d = 10 = \mathbf{2+2}+\mathbf{3+3}$: both pairs annihilate $\to$ nothing.
\end{example}

\subsection{Swap elimination: Topological proof}

\begin{theorem}[Grassmannian Attractor]
\label{thm:grassmann}
If $a = b$, the Grassmannian $\text{Gr}(a, 2a)$ of causal decompositions has a $\mathbb{Z}_2$-invariant fixed locus that is a topological attractor, trapping the dynamics in a structureless configuration.
\end{theorem}

\begin{proof}
\textbf{Step 1.}
The space of decompositions is $\mathcal{G} = \text{Gr}(a, d) = \text{U}(d)/(\text{U}(a) \times \text{U}(b))$.

\textbf{Step 2.}
When $a = b$: the complement map $\tau: V_S \mapsto V_S^\perp$ is a well-defined involution on $\mathcal{G}$ (since $\dim V_S^\perp = b = a$). For $a \neq b$: $\tau$ maps $\text{Gr}(a, d) \to \text{Gr}(b, d) \neq \text{Gr}(a, d)$. No involution.

\textbf{Step 3.}
The fixed locus $\mathcal{G}^\tau = \{V_S : V_S = V_S^\perp\} = \text{U}(a)/\text{O}(a)$ is non-empty for all $a \geq 1$.

\textbf{Step 4.}
Entropy maximization drives the system toward $\mathcal{G}^\tau$ (the maximally disordered symmetric configuration). Once there, $V_S \cong V_T$, freezing is symmetric, $\hbar$ is constant. No physics.

For $a \neq b$: no involution $\Rightarrow$ no attractor $\Rightarrow$ asymmetric dynamics is forced.
\end{proof}

\subsection{Gravity requires $d \geq 5$}

\begin{proposition}
For $d_\text{spacetime} = d - 1 \leq 3$, gravity has zero propagating degrees of freedom.
\end{proposition}

\begin{proof}
Graviton polarizations: $n_\text{grav} = d_\text{spacetime}(d_\text{spacetime} - 3)/2$. For $d_\text{spacetime} \leq 3$: $n_\text{grav} \leq 0$. No gravitational waves, no black hole horizons. The Gram determinant $\det(G_h)$ encodes no dynamical information: spacetime curvature is fully determined by local matter content.
\end{proof}

\subsection{The uniqueness theorem}

\begin{theorem}[Uniqueness of $d = 5$]
\label{thm:d5}
The unique integer $d$ satisfying:
\begin{enumerate}[label=(\roman*)]
\item $d \geq 5$ (propagating gravity),
\item $d = a + b$ with $a, b \geq 2$ both atomic and $a \neq b$,
\end{enumerate}
is $d = 5$, with $(a, b) = (3, 2)$.
\end{theorem}

\begin{proof}
Atoms $= \{2, 3\}$ (Lemma~\ref{lem:atoms}). Distinct pairs $= \{(2, 3)\}$. Sum $= 5 \geq 5$.\;\;$\square$
\end{proof}

\subsection{The cosmic attractor}

The uniqueness theorem has a dynamical interpretation. Starting from $\text{SU}(N)$ for any $N$:
\begin{equation}
N = \underbrace{2 + 2 + \cdots + 2}_{a\text{ pairs}} + \underbrace{3 + 3 + \cdots + 3}_{b\text{ pairs}} + \text{remainder}
\end{equation}

All repeated atoms annihilate by Corollary~\ref{cor:annihilate}. At most one 2 and one 3 survive:
\begin{equation}
\text{SU}(N) \;\xrightarrow{\;\text{swap annihilation}\;}\; \text{SU}(3) \times \text{SU}(2) \times \text{U}(1)
\end{equation}
\emph{independent of $N$}. The Standard Model gauge group is the universal attractor of dimension reduction.

\subsection{The U(1) glue}
\label{sec:glue}

After the split $\mathbb{C}^5 = \mathbb{C}^3 \oplus \mathbb{C}^2$, the U(1) factor is the relative phase between the two sectors --- the unique element of $\mathbb{C}$ that Chapter~\ref{ch:whyC} established. The tracelessness of hypercharge:
\begin{equation}
n_S \cdot y_\text{color} + n_T \cdot y_\text{weak} = 0 \qquad\Rightarrow\qquad \frac{y_\text{color}}{y_\text{weak}} = -\frac{n_T}{n_S} = -\frac{2}{3}
\end{equation}
enforces exact charge conservation: what one sector gives, the other takes. The electromagnetic force is not an accident but the \emph{necessary glue} connecting space and time.

\subsection{Summary}

\begin{center}
\begin{tabular}{cccl}
\hline
$d$ & Decomposition & Physics? & Reason \\
\hline
$\leq 4$ & --- & No & Gravity: 0 degrees of freedom \\
4 & $2+2$ & No & Swap $\to$ symmetric $\to$ trivial \\
\textbf{5} & $\mathbf{2+3}$ & \textbf{Yes} & \textbf{Unique distinct atomic pair} \\
6 & $3+3$ & No & Swap $\to$ trivial \\
7 & $2+2+3$ & No & Repeated 2's $\to$ annihilate \\
$\geq 6$ & Any & No & All repeated atoms annihilate $\to$ only $(3,2)$ \\
$N$ & Any & $\to \mathbf{5}$ & Universal attractor \\
\hline
\end{tabular}
\end{center}

\noindent The spacetime dimension $d_\text{spacetime} = 4$ is a theorem, not a parameter.
